\chapter{श्रेणी और समांतर परिपथ}\label{ch:g5:series-parallel-circuits}

एक \omterm{def:g3:first-electric-circuit:bulb}{बल्ब} मरा और सजावटी बत्तियों की पूरी लड़ी अँधेरी हो गई --- अकेले
फंदे का शाप। फिर भी तुम्हारी रसोई का कोई \omterm{def:g3:first-electric-circuit:bulb}{बल्ब} फुँक जाए तो बाक़ी पूरा
घर जगमगाता रहता है। किसी ने तुम्हारा घर ज़्यादा चतुराई से जोड़ा है।
आज हम वही तरीक़ा सीखेंगे: \omterm{def:g5:series-parallel-circuits:parallel}{शाखाओं} वाला रास्ता।

\section{जोड़ने के दो तरीक़े}

पहला तरीक़ा तुम जानते हो: \emph{\omterm{def:g4:series-circuits:series}{श्रेणी में}}, जहाँ सारे खिलाड़ी एक ही
फंदे पर पिरोए रहते हैं, एक ही माला के मनकों की तरह। अब दूसरा।

\begin{definition}[समांतर में]\label{def:g5:series-parallel-circuits:parallel}
खिलाड़ी \emph{समांतर में}\index{समांतर परिपथ} तब होते हैं जब परिपथ
अलग-अलग \emph{शाखाओं}\index{शाखा} में बँट जाता है --- हर शाखा पर एक
खिलाड़ी --- जो अलग होकर फिर से मिल जाती हैं। जिन बिंदुओं पर सड़कें
फूटती या फिर से मिलती हैं उन्हें \emph{संधि}\index{संधि} कहते हैं। हर
शाखा \omterm{def:g3:first-electric-circuit:battery}{बैटरी} के एक सिरे से दूसरे सिरे तक जाती एक पूरी सड़क होती है।
\end{definition}

\begin{omfigure}
\begin{tikzpicture}[scale=0.8]
  \draw (0,0) to[battery1, l=बैटरी] (0,3.2)
    to[short] (2.2,3.2);
  \draw (2.2,3.2) to[short] (2.2,2.4)
    to[lamp, l=बल्ब 1] (5.4,2.4) to[short] (5.4,3.2);
  \draw (2.2,3.2) to[short] (2.2,4.0)
    to[lamp, l=बल्ब 2] (5.4,4.0) to[short] (5.4,3.2);
  \draw (5.4,3.2) to[short] (7.0,3.2) to[short] (7.0,0)
    to[short] (0,0);
  \fill (2.2,3.2) circle (2pt);
  \fill (5.4,3.2) circle (2pt);
  \node[above, font=\small] at (3.6,4.95) {दो शाखाएँ};
\end{tikzpicture}

{\small \omterm{def:g5:series-parallel-circuits:parallel}{समांतर में} दो \omterm{def:g3:first-electric-circuit:bulb}{बल्ब}: सड़क एक \omterm{def:g5:series-parallel-circuits:parallel}{संधि} पर फूटती है, हर \omterm{def:g3:first-electric-circuit:bulb}{बल्ब} अपनी
\omterm{def:g5:series-parallel-circuits:parallel}{शाखा} पर सवार होता है, और \omterm{def:g3:first-electric-circuit:battery}{बैटरी} तक लौटने से पहले शाखाएँ फिर मिल जाती
हैं।}
\end{omfigure}

\begin{method}[इसे बनाना]\label{met:g5:series-parallel-circuits:build}
एक \omterm{def:g3:first-electric-circuit:battery}{बैटरी}, दो \omterm{def:g3:first-electric-circuit:bulb}{बल्ब} और तारों से:
\begin{enumerate}
  \item पहले \omterm{def:g3:first-electric-circuit:bulb}{बल्ब} को अकेले, सीधे फंदे की तरह, \omterm{def:g3:first-electric-circuit:battery}{बैटरी} से जोड़ो ---
        यह \omterm{def:g5:series-parallel-circuits:parallel}{शाखा} एक हुई;
  \item अब दूसरे \omterm{def:g3:first-electric-circuit:bulb}{बल्ब} को \emph{सीधे पहले \omterm{def:g3:first-electric-circuit:bulb}{बल्ब} के दोनों सिरों पर}
        जोड़ो: \omterm{def:g3:first-electric-circuit:bulb}{बल्ब} 1 की हर ओर से एक तार \omterm{def:g3:first-electric-circuit:bulb}{बल्ब} 2 की मेल खाती ओर तक
        --- यह \omterm{def:g5:series-parallel-circuits:parallel}{शाखा} दो हुई;
  \item हर \omterm{def:g5:series-parallel-circuits:parallel}{शाखा} का उँगली से पीछा करो: $+$ से चलकर हर सड़क को $-$ तक
        पहुँचने से पहले ठीक एक \omterm{def:g3:first-electric-circuit:bulb}{बल्ब} पार करना चाहिए।
\end{enumerate}
दोनों \omterm{def:g3:first-electric-circuit:bulb}{बल्ब} जलते हैं --- और दोनों उतने ही चमकीले, जितना अकेला \omterm{def:g3:first-electric-circuit:bulb}{बल्ब}
चमकता।
\end{method}

\section{शाखाओं के नियम}

\begin{proposition}[समांतर परिपथ के नियम]\label{prop:g5:series-parallel-circuits:rules}
जब हर खिलाड़ी अपनी \omterm{def:g5:series-parallel-circuits:parallel}{शाखा} पर हो:
\begin{enumerate}
  \item \emph{स्वतंत्रता}: किसी एक \omterm{def:g5:series-parallel-circuits:parallel}{शाखा} की दरार सिर्फ़ उसी \omterm{def:g5:series-parallel-circuits:parallel}{शाखा} को
        रोकती है --- एक \omterm{def:g3:first-electric-circuit:bulb}{बल्ब} खोल दो, दूसरा जलता रहता है;
  \item \emph{पूरी चमक}: हर \omterm{def:g3:first-electric-circuit:bulb}{बल्ब} ऐसे चमकता है मानो वह अकेला हो ---
        \omterm{def:g5:series-parallel-circuits:parallel}{शाखाओं} के बीच \omterm{def:g3:first-electric-circuit:battery}{बैटरी} का धक्का बँटता नहीं;
  \item \emph{हुक्म की चौकियाँ}: किसी \omterm{def:g5:series-parallel-circuits:parallel}{शाखा} के भीतर लगा \omterm{ex:g3:first-electric-circuit:switch}{स्विच} सिर्फ़
        उसी \omterm{def:g5:series-parallel-circuits:parallel}{शाखा} को हुक्म देता है; और साझा सड़क पर, फूट से पहले लगा
        \omterm{ex:g3:first-electric-circuit:switch}{स्विच} सबको। (क़ीमत: \omterm{def:g3:first-electric-circuit:battery}{बैटरी} हर \omterm{def:g5:series-parallel-circuits:parallel}{शाखा} को एक साथ खिलाती है,
        इसलिए वह ज़्यादा तेज़ी से चुकती है।)
\end{enumerate}
\end{proposition}

\begin{example}[खोलकर देखने वाली जाँच, दोबारा]\label{ex:g5:series-parallel-circuits:test}
समांतर वाली जोड़ी के साथ ठीक वही करो जिसने कभी श्रेणी वाली जोड़ी की
जान ले ली थी: \omterm{def:g3:first-electric-circuit:bulb}{बल्ब} 2 खोल दो। \omterm{def:g3:first-electric-circuit:bulb}{बल्ब} 1 पलक तक नहीं झपकाता --- उसकी अपनी
\omterm{def:g5:series-parallel-circuits:parallel}{शाखा}, यानी उसकी अपनी पूरी सड़क, छुई ही नहीं गई। उसे वापस कस दो और
इस बार \omterm{def:g3:first-electric-circuit:bulb}{बल्ब} 1 खोलो: \omterm{def:g3:first-electric-circuit:bulb}{बल्ब} 2 शांति से जलता रहता है। अकेले फंदे का शाप
उतर गया: \omterm{def:g5:series-parallel-circuits:parallel}{समांतर में} हर खिलाड़ी अपनी \omterm{def:g5:series-parallel-circuits:parallel}{शाखा} से मतलब रखता है।
\end{example}

\begin{omfigure}
\begin{tikzpicture}[scale=0.8]
  \draw (0,0) to[battery1] (0,3.2)
    to[cspst, l=मुख्य स्विच] (2.6,3.2);
  \draw (2.6,3.2) to[short] (2.6,2.4)
    to[lamp, l=बल्ब 1] (5.8,2.4) to[short] (5.8,3.2);
  \draw (2.6,3.2) to[short] (2.6,4.0)
    to[spst, l=शाखा का स्विच] (4.2,4.0)
    to[lamp] (5.8,4.0) to[short] (5.8,3.2);
  \draw (5.8,3.2) to[short] (7.4,3.2) to[short] (7.4,0)
    to[short] (0,0);
  \fill (2.6,3.2) circle (2pt);
  \fill (5.8,3.2) circle (2pt);
\end{tikzpicture}

{\small हुक्म की चौकियाँ: फूट से पहले लगा मुख्य \omterm{ex:g3:first-electric-circuit:switch}{स्विच} दोनों \omterm{def:g5:series-parallel-circuits:parallel}{शाखाओं} पर
राज करता है; \omterm{def:g5:series-parallel-circuits:parallel}{शाखा} का \omterm{ex:g3:first-electric-circuit:switch}{स्विच} सिर्फ़ अपने \omterm{def:g3:first-electric-circuit:bulb}{बल्ब} पर --- यहाँ वह खुला है,
इसलिए \omterm{def:g3:first-electric-circuit:bulb}{बल्ब} 1 अकेला चमक रहा है।}
\end{omfigure}

\begin{example}[तुम्हारा घर एक समांतर परिपथ है]\label{ex:g5:series-parallel-circuits:house}
तुम्हारे घर का हर लैंप, हर सॉकेट और हर उपकरण एक विशाल \omterm{def:g5:series-parallel-circuits:parallel}{समांतर परिपथ} की
अपनी-अपनी \omterm{def:g5:series-parallel-circuits:parallel}{शाखा} पर सवार है। इसीलिए रसोई का \omterm{def:g3:first-electric-circuit:bulb}{बल्ब} फुँक जाने पर भी बैठक
बिना रुके पढ़ती रहती है; इसीलिए हर कमरे का \omterm{ex:g3:first-electric-circuit:switch}{स्विच} सिर्फ़ अपनी बत्ती को
हुक्म देता है; और इसीलिए \omterm{def:g2:measuring-length:units}{मीटर} के पास फ़्यूज़-पेटी में लगा एक मुख्य
\omterm{ex:g3:first-electric-circuit:switch}{स्विच} मरम्मत के लिए पूरे घर की बिजली एक साथ काट सकता है। आजकल की
सजावटी बत्तियाँ भी इसी क्लब में शामिल हो गई हैं: उनके \omterm{def:g3:first-electric-circuit:bulb}{बल्ब} इस तरह जुड़े
होते हैं कि एक \omterm{def:g3:first-electric-circuit:bulb}{बल्ब} मर जाने से पूरी लड़ी अँधेरी नहीं होती।
\end{example}

\begin{remark}[श्रेणी और समांतर, आमने-सामने]\label{rem:g5:series-parallel-circuits:compare}
एक दरार: श्रेणी --- सब अँधेरा; समांतर --- सिर्फ़ एक \omterm{def:g5:series-parallel-circuits:parallel}{शाखा} अँधेरी।
ज़्यादा \omterm{def:g3:first-electric-circuit:bulb}{बल्ब}: श्रेणी --- सब मंद; समांतर --- सब पूरी चमक पर, पर \omterm{def:g3:first-electric-circuit:battery}{बैटरी}
ज़्यादा तेज़ी से चुकती है। सबके लिए एक \omterm{ex:g3:first-electric-circuit:switch}{स्विच}: श्रेणी --- कहीं भी;
समांतर --- सिर्फ़ फूट से पहले। इनमें से कोई तरीक़ा ``बेहतर'' नहीं है:
टॉर्च ख़ुशी-ख़ुशी \omterm{def:g4:series-circuits:series}{श्रेणी में} है (एक \omterm{def:g3:first-electric-circuit:bulb}{बल्ब}, एक \omterm{ex:g3:first-electric-circuit:switch}{स्विच}), और घर मजबूरन
\omterm{def:g5:series-parallel-circuits:parallel}{समांतर में}। किसी भी परिपथ के बारे में बिजली वाले का पहला सवाल बस इतना
होता है: \emph{सड़क फूटती कहाँ है?}
\end{remark}

\section{परिपथ पढ़ना}

\begin{method}[सड़कों का पीछा करना]\label{met:g5:series-parallel-circuits:follow}
किसी परिपथ-चित्र में किसी \omterm{def:g3:first-electric-circuit:bulb}{बल्ब} का भविष्य बताने के लिए:
\begin{enumerate}
  \item उँगली \omterm{def:g3:first-electric-circuit:battery}{बैटरी} के $+$ पर रखो और $-$ तक जाने वाली हर संभव सड़क
        चलो, और हर सड़क पर मिलने वाले खिलाड़ी गिनते जाओ;
  \item कोई \omterm{def:g3:first-electric-circuit:bulb}{बल्ब} तब जलता है जब कम से कम एक अटूट सड़क उसमें से गुज़रती
        हो;
  \item जो खिलाड़ी हर सड़क साझा करते हैं वे आपस में \omterm{def:g4:series-circuits:series}{श्रेणी में} हैं;
        और अलग-अलग फूटों पर बैठे खिलाड़ी \omterm{def:g5:series-parallel-circuits:parallel}{समांतर में}।
\end{enumerate}
भविष्य बताने से पहले सड़कें चल लो --- चित्र की शक्ल देखकर अंदाज़ा
लगाना ही वह तरीक़ा है जिससे परिपथ लोगों को धोखा देते हैं।
\end{method}

\section{अभ्यास}

\begin{exercise}[$\star$]\label{exo:g5:series-parallel-circuits:1}
\omterm{def:g5:series-parallel-circuits:parallel}{शाखा} क्या है? \omterm{def:g5:series-parallel-circuits:parallel}{संधि} क्या है? दो बल्बों वाले \omterm{def:g5:series-parallel-circuits:parallel}{समांतर परिपथ} में कितनी
\omterm{def:g5:series-parallel-circuits:parallel}{शाखाओं} पर \omterm{def:g3:first-electric-circuit:bulb}{बल्ब} होते हैं?
\end{exercise}

\begin{exercise}[$\star$]\label{exo:g5:series-parallel-circuits:2}
\omterm{def:g5:series-parallel-circuits:parallel}{समांतर परिपथ} के तीनों नियम बताओ।
\end{exercise}

\begin{exercise}[$\star$]\label{exo:g5:series-parallel-circuits:3}
\omterm{def:g5:series-parallel-circuits:parallel}{समांतर में} लगे दो \omterm{def:g3:first-electric-circuit:bulb}{बल्ब}: \omterm{def:g3:first-electric-circuit:bulb}{बल्ब} 1 खोल दिया जाता है। \omterm{def:g3:first-electric-circuit:bulb}{बल्ब} 2 क्या करता
है? और श्रेणी वाली जोड़ी में उसने वहाँ क्या किया था, और यह फ़र्क़
क्यों?
\end{exercise}

\begin{exercise}[$\star$]\label{exo:g5:series-parallel-circuits:4}
एक जैसे दो \omterm{def:g3:first-electric-circuit:bulb}{बल्ब}: जोड़ी क \omterm{def:g4:series-circuits:series}{श्रेणी में} है, जोड़ी ख \omterm{def:g5:series-parallel-circuits:parallel}{समांतर में}, और
बैटरियाँ एक जैसी। हर \omterm{def:g3:first-electric-circuit:bulb}{बल्ब} के हिसाब से कौन-सी जोड़ी ज़्यादा चमकती है?
और कौन-सी अपनी \omterm{def:g3:first-electric-circuit:battery}{बैटरी} ज़्यादा तेज़ी से चुकाती है?
\end{exercise}

\begin{exercise}[$\star$]\label{exo:g5:series-parallel-circuits:5}
दो स्विचों वाले चित्र में हर \omterm{def:g3:first-electric-circuit:bulb}{बल्ब} क्या करता है जब: दोनों \omterm{ex:g3:first-electric-circuit:switch}{स्विच} बंद
हों? \ मुख्य खुला हो? \ मुख्य बंद और \omterm{def:g5:series-parallel-circuits:parallel}{शाखा} का \omterm{ex:g3:first-electric-circuit:switch}{स्विच} खुला हो?
\end{exercise}

\begin{exercise}[$\star$]\label{exo:g5:series-parallel-circuits:6}
घर \omterm{def:g4:series-circuits:series}{श्रेणी में} नहीं, \omterm{def:g5:series-parallel-circuits:parallel}{समांतर में} क्यों जोड़ा जाता है? इस पर टिकी दो
रोज़मर्रा की सुविधाएँ बताओ।
\end{exercise}

\begin{exercise}[$\star$]\label{exo:g5:series-parallel-circuits:7}
घर के परिपथ में \omterm{ex:g3:first-electric-circuit:switch}{स्विच} कहाँ लगा हो तो वह सब कुछ एक साथ सँभाल लेता है?
उस \omterm{ex:g3:first-electric-circuit:switch}{स्विच} का रोज़ का नाम और काम क्या है?
\end{exercise}

\begin{exercise}[$\star$]\label{exo:g5:series-parallel-circuits:8}
सिरहाने का लैंप और उसका दीवार वाला सॉकेट: वे छत की बत्ती के साथ
\omterm{def:g4:series-circuits:series}{श्रेणी में} हैं या \omterm{def:g5:series-parallel-circuits:parallel}{समांतर में}? रोज़ के अनुभव से तुम्हें कैसे पता?
\end{exercise}

\begin{exercise}[$\star\star$]\label{exo:g5:series-parallel-circuits:9}
संकेतों से किसी गलियारे का परिपथ बनाओ: एक \omterm{def:g3:first-electric-circuit:battery}{बैटरी}, एक \omterm{def:g3:first-electric-circuit:bulb}{बल्ब}, और यह
सुविधा --- \omterm{def:g3:first-electric-circuit:bulb}{बल्ब} गलियारे के \emph{किसी भी} सिरे पर लगे \omterm{ex:g3:first-electric-circuit:switch}{स्विच} से जलना
और बुझना चाहिए\dots\ ईमानदारी से कोशिश करो; अगर तुम्हारे डिज़ाइन बार-बार
नाकाम हों, तो बताओ कि कठिन क्या है। (सचमुच की दो-तरफ़ा स्विचिंग के
लिए हमारे \omterm{ex:g3:first-electric-circuit:switch}{स्विच} से ज़्यादा चतुर \omterm{ex:g3:first-electric-circuit:switch}{स्विच} चाहिए --- और यहाँ कोशिश उत्तर से
ज़्यादा सिखाती है।)
\end{exercise}

\begin{exercise}[$\star\star$]\label{exo:g5:series-parallel-circuits:10}
आजकल की सजावटी बत्तियों की लड़ी एक \omterm{def:g3:first-electric-circuit:bulb}{बल्ब} मर जाने पर भी जलती रहती है
--- पर डिब्बे पर चेतावनी लिखी है: ``तीन से ज़्यादा \omterm{def:g3:first-electric-circuit:bulb}{बल्ब} निकालकर लड़ी
मत चलाइए।'' \omterm{def:g5:series-parallel-circuits:parallel}{शाखाओं} के नियमों से यह सुविधा समझाओ, और उससे कठिन काम यह
कि उस चेतावनी की कोई समझदार वजह भी बताओ। (सोचो कि हर बची हुई \omterm{def:g5:series-parallel-circuits:parallel}{शाखा} को
\omterm{def:g3:first-electric-circuit:battery}{बैटरी} से क्या मिलता है, और \omterm{def:g3:first-electric-circuit:battery}{बैटरी} को ख़ुद कितना ज़ोर लगाना पड़ता है।)
\end{exercise}

\begin{exercise}[$\star\star$]\label{exo:g5:series-parallel-circuits:11}
दो स्विचों वाले चित्र में एक बिजली वाला ऐसा एक और \omterm{def:g3:first-electric-circuit:bulb}{बल्ब} जोड़ना चाहता है
जो \emph{सिर्फ़} तब जले जब \omterm{def:g5:series-parallel-circuits:parallel}{शाखा} का \omterm{ex:g3:first-electric-circuit:switch}{स्विच} बंद हो, और बाक़ी समय अँधेरा
रहे --- मुख्य \omterm{ex:g3:first-electric-circuit:switch}{स्विच} बंद होने पर भी। नया \omterm{def:g3:first-electric-circuit:bulb}{बल्ब} सड़क के किस हिस्से पर
जोड़ना होगा, और तब वह \omterm{def:g3:first-electric-circuit:bulb}{बल्ब} 2 के साथ \omterm{def:g4:series-circuits:series}{श्रेणी में} होगा या \omterm{def:g5:series-parallel-circuits:parallel}{समांतर में}?
\end{exercise}
